\documentclass[12pt,a4paper]{article}
%\usepackage[UTF8]{ctex}
\usepackage{amsfonts,amsmath,amsthm,amssymb}
\usepackage{sfmath}
\usepackage{hyperref}
\renewcommand{\familydefault}{\sfdefault}
\usepackage{geometry}
\geometry{left=0cm,right=0cm,top=0cm,bottom=0.1cm, footskip=0cm, includefoot=true}
%\usepackage[contents = \thepage, color=yellow, angle=0, scale = 20]{background}
\pagestyle{empty}
%\usepackage{fancyhdr}
%\pagestyle{fancy}
%\fancyhead{}
%\fancyfoot{}
%\cfoot{}
%\lfoot{{\small \color{red} \slshape \thepage 页}}
%\fancyfoot[R]{\thepage}

\usepackage[all]{xy}
\usepackage{color}

\renewenvironment{proof}{{\noindent\it\textbf{Proof}}\quad\it}{\hfill$\square$}

%\theoremstyle{definition}

\newtheorem{thm}[equation]{Theorem}
%\newtheorem{theorem}[equation]{Theorem}
\newtheorem{cor}[equation]{Corollary}
\newtheorem{lem}[equation]{Lemma}
\newtheorem{prop}[equation]{Proposition}
\newtheorem{defn}[equation]{Definition}
\newtheorem{rem}[equation]{Remark}
\newtheorem{conj}[equation]{Conjecture}

\newcommand{\Z}{\mathbb{Z}}

\title{Unboundedness}
\author{}


\begin{document}

\sffamily
\maketitle
\thispagestyle{empty}

\abstract{}
\tableofcontents

\large

\begin{sloppy}
\clubpenalty=-100
\widowpenalty=-100

\section{stabilize}

\begin{theorem}
\begin{thm}
Let $E_1, A_1$ and $E_2, A_2$ be converging spectral sequences, and $f$ a morphism.

Let $s_1<s_2$ two filtration indexes.

Then if $s_1> s +r+1$, the filtration $s$ commponent of the ess for $s_1$-truncation is isomorphism to the ess for $s_2$-truncation.
\end{thm}
\begin{proof}
    By definition, in ess,
    $V$ is the associated graded of $A_i$, iin filtration degree $s$, this is the same as the graded piece of the truncated ones.

    $Z_r$ is defined to be the projection of the kernel of 
    $$d: F^sA_1 \rightarrow F^sA_1/F^{s+r-1}A_1$$
if $s+r<s_1$, this definition is not affected by truncation.

$B_r$ is defined to be the image of 
$$d: Z_{r-1}\cap Fil^{s-r+1}A_1 \rightarrow Fil^sA_2\rightarrow E_r$$

This also not affected by truncation

$d_r$ is define by 
$$d: Fil^sA_1\rightarrow Fil^{s+r}A_2$$
this is not affected by the truncation.

\end{proof}



\begin{thebibliography}{99}


\end{thebibliography}

\end{sloppy}
\end{document}
